\section{Introduction}
\label{sec:intro}

A minimal hypersurface is a critical point of the area functional
in codimension one. On a closed Riemannian manifold
$(M^{n+1}, g)$ with $H_n(M;\mathbb{Z}) = 0$ (for instance
$M = S^{n+1}$), direct minimization of the $n$-dimensional
Hausdorff measure over integral currents yields only the empty
hypersurface; a non-trivial minimal hypersurface must be
constructed as a saddle-type critical point of the area
functional.
\me{Almgren--Pitts~\cite{Pit81} give a \emph{min-max} scheme on
varifolds; regularity is established by Schoen--Simon~\cite{SS81}.}

\begin{theorem}[Almgren--Pitts min-max theorem~\cite{Pit81,SS81}]
\label{thm:minmax}
Every closed Riemannian manifold $(M^{n+1}, g)$ admits a
nontrivial embedded minimal hypersurface $\Sigma \subset M$ whose
singular set $\mathrm{Sing}\,\Sigma$ has Hausdorff dimension at
most $n - 7$.
\end{theorem}

A key step in Pitts's proof, recurring in
\cite{CDL03,DLT13,DLR17,DL12,MN14}, is the Almgren--Pitts
combinatorial argument, stated in the Colding--De~Lellis
formulation as follows.

\begin{theorem}[Almgren--Pitts combinatorial Argument,
{\cite[Prop.~3.4]{CDL03}}]
\label{thm:AP-combinatorial}
Let $\Lambda$ be a homotopically closed family of sweepouts of
$M^{n+1}$, and $m_0 = \inf_\Lambda \mathcal{F}$. There exists a
min-max sequence $\{\Gamma^N\}$ converging as varifolds to a
stationary integer rectifiable varifold of mass $m_0$, with
$\Gamma^N$
$1/N$-almost-minimizing in every pair
$(U^1, U^2) \in \mathcal{CO}$ of $\mathcal{H}^n$-disjoint open
sets, for all $N$ sufficiently large.
\end{theorem}

\me{The proof of Theorem~\ref{thm:AP-combinatorial} reduces,
following~\cite[\S3.2]{DLT13}, to the following claim on the
sweepout:\footnote{Throughout this paper, ``sweepout'' denotes a
one-parameter family $\{\Gamma_t\}_{t \in [0,1]}$ in the sense of
\cite[Def.~3.1]{DLT13}; the multi-parameter min-max setting of
\cite{MN14} requires separate combinatorial machinery and lies
outside our scope.}}

\begin{quote}
\textbf{Claim.} For each $N$ sufficiently large, there exists
$t_N \in [0,1]$ such that $\Gamma^N := \Gamma^N_{t_N}$ is
$1/N$-almost-minimizing in every $(U^1, U^2) \in \mathcal{CO}$
and $\mathcal{H}^n(\Gamma^N) \ge m_0 - 1/N$.
\end{quote}

The Claim is proved by contradiction: from its negation, every
$t$ with $\mathcal{H}^n(\partial \Omega^N_t) \ge m_0 - 1/N$
admits a \emph{local replacement} with saving $1/(4N)$
\cite[Lemma~3.1]{DLT13}; a combinatorial cover-refinement and
bookkeeping argument on $[0,1]$ then assembles from these local
replacements a modified sweepout with supremum strictly below
$m_0$,
contradicting $m_0 = \inf_\Lambda \mathcal{F}$. The combinatorial
argument depends on no geometric measure theory, no admissible
class, and no ambient regularity; it is a bookkeeping
construction on continuous real-valued functions, yet has no
canonical reference.

The combinatorial core of the proof of the Claim, which is the
precise theorem we formalize, reads as follows.

\begin{theorem}[Combinatorial core]
\label{thm:intro-scalar}
Let $f : [0,1] \to \mathbb{R}$ be continuous with
$m_0 = \sup f$, and let $N \in \mathbb{N}$ with $N > 0$. Suppose
that at every $t \in [0,1]$ with $f(t) \ge m_0 - 1/N$ there is a
\emph{local witness}: an open neighborhood $U_t \ni t$ and a
continuous replacement energy $E_t : [0,1] \to \mathbb{R}$
satisfying $f(s) - E_t(s) \ge 1/(4N)$ for every $s \in U_t$.
Then there exist a finite family of closed pieces
$\{p_l \subseteq [0,1]\}_{l \in \iota}$, continuous replacement
energies $\{E_l : [0,1] \to \mathbb{R}\}_{l \in \iota}$, and
savings $\{s_l \ge 1/(4N)\}_{l \in \iota}$ satisfying:
\begin{enumerate}
\item[\textup{(I)}] $f - E_l \ge s_l$ on $p_l$;
\item[\textup{(II)}] every $t \in [0,1]$ belongs to at most
  two of the pieces $p_l$;
\item[\textup{(III)}] $\{t \in [0,1] : f(t) \ge m_0 - 1/N\}
  \subseteq \bigcup_l p_l$.
\end{enumerate}
\end{theorem}

\begin{remark}
\label{rem:intro-core-nontrivial}
\me{A natural compactness sketch of Theorem~\ref{thm:intro-scalar}
runs: $f^{-1}([m_0 - 1/N, \infty))$ is compact, so finitely many
witness neighborhoods cover it; pass to an irredundant subfamily
and take $p_l := \overline{U_{t_i}}$. Conditions~(I) and~(III) hold
immediately; (II) does not. Even in dimension one, sub-selection
alone need not deliver multiplicity~$\le 2$ on \emph{closed}
pieces: $I_1 = (0, 0.4)$, $I_2 = (0.3, 0.5)$, $I_3 = (0.4, 0.7)$
together cover $[0.05, 0.65]$, every member is needed, and the
open multiplicity is $\le 2$, yet $0.4 \in
\overline{I_1} \cap \overline{I_2} \cap \overline{I_3}$. The
witness hypothesis offers no leverage to avoid such configurations:
nothing constrains the $U_t$ to be spaced.}

The non-trivial content of Theorem~\ref{thm:intro-scalar} is the
\emph{construction} of a cover $\{p_l\}$ \me{which precludes
them}. A Lebesgue-number argument on the compact near-critical
set produces an initial cover by intervals $\{I_i\}$ of controlled
radius satisfying the skip-$2$ spacing condition~(a)
of~\cite{DLT13}, \me{$\overline{I_i} \cap \overline{I_{i+2}}
= \emptyset$ --- exactly ruling out the boundary collision above}.
A bounded smallest-index induction refines $\{I_i\}$ to pieces
$\{J_k\}$, where each new piece is either $I_{i^*}$ or
$I_{i^*} \setminus I_{\sigma(\ell)}$; \me{the set-difference branch
is what makes this strictly more than sub-selection}. A parity
argument on~(a) then transports the skip-$2$ disjoint-closure
property to the two-fold overlap~(II) on $\{\overline{J_k}\}$;
continuity extends the per-piece saving from $J_k$ to
$\overline{J_k}$. The bookkeeping corollary below (the existence
of a competitor $f'$ with reduced supremum) is a three-line
arithmetic consequence.
\end{remark}

\begin{remark}[Why the construction, not just the conclusion]
\label{rem:why-construction}
\me{At the abstract scalar level alone, Theorem~\ref{thm:intro-scalar}
admits a strictly shorter proof}: a finite subcover by witness
intervals, refined to a partition by collecting all interval
endpoints, yields closed pieces of multiplicity~$\le 2$ (sharing
only endpoints) carrying the witness data. \me{Such a partition
suffices for the scalar conclusion because the multiplicity bound
(II) is structurally recorded but \emph{not arithmetically
load-bearing} for $\sup f' \le m_0 - 1/(4N)$ --- any cover with
per-piece saving $\ge 1/(4N)$ delivers (c) of
Corollary~\ref{thm:intro-core} (\texttt{Cover.lean} states this
explicitly).}

\me{What a partition discards is precisely what the formalization
is meant to extract.} The Almgren--Pitts combinatorial argument
as it appears in \cite{CDL03,DLT13,DLR17,DL12,MN14} is not the
scalar conclusion of Theorem~\ref{thm:intro-scalar} but the
\emph{construction} that proves it: overlapping intervals
$\{I_i\}$ with spacing condition~(a), an injective refinement
$\sigma$, and the parity rescue of two-fold overlap on closures.
\me{The geometric lift in the ambient GMT setting --- DLT's
modified sweepout $\{\partial \Omega'_t\}$, defined by blending
two replacements on $J_i \cap J_{i+1}$ --- requires
\emph{positive-measure} overlap and breaks at a partition's
measure-zero endpoint sharing} (the blend
$\Omega'_t = [\Omega_t \setminus (U_i' \cup U_{i+1}')] \cup
[\Omega_{i,t} \cap U_i'] \cup [\Omega_{i+1,t} \cap U_{i+1}']$
needs both replacements simultaneously available on a region of
positive $t$-measure to remain $t$-continuous).

\me{The library therefore ships a two-tier architecture: the
DLT-faithful path under \texttt{CombArg.OneDim} produces a
structured \texttt{DLTCover} that exposes the spacing /
refinement / $\sigma$-injectivity data a geometric consumer
needs, with the abstract scalar theorem obtained as a
downgrade; the partition shortcut is also formalized, as
\texttt{CombArg.Scalar.exists\_refinement\_partition},
specifically as evidence that the DLT machinery is dispensable
at the abstract scalar level alone --- a dependency-graph fact
verified by \texttt{lake build}: the partition module imports
neither \texttt{OneDim/SpacedIntervals} nor
\texttt{OneDim/Induction}.}
\end{remark}

Translation from the De~Lellis--Tasnady setting to
Theorem~\ref{thm:intro-scalar}: the boundary $\partial \Omega_t$
becomes a continuous $f : [0,1] \to \mathbb{R}$; the
Hausdorff-measure hypothesis becomes $m_0 = \sup f$; the nested
pair $(A_t, B_t, U_t')$ collapses to a single neighborhood $U_t$;
the replacement family $\{\partial \Omega_{t,\tau}\}$ collapses
to a single continuous $E_t$; and the almost-minimizing
hypothesis, which in the original proof is the input to
Lemma~3.1 of~\cite{DLT13}, is absorbed into the existence of the
local witness (the caller must supply witnesses wherever they
would be produced in the ambient GMT setting). The output of
Theorem~\ref{thm:intro-scalar} is the combinatorially assembled
cover data $\{(p_l, E_l, s_l)\}$; the modified sweepout in the
original proof is the geometric lift of the scalar competitor
produced by the corollary below.

\begin{corollary}[Sup reduction]
\label{thm:intro-core}
Under the hypotheses of Theorem~\ref{thm:intro-scalar}, with the
pieces $\{p_l\}$ produced there, there exists
$f' : [0,1] \to \mathbb{R}$ satisfying:
\begin{enumerate}
\item[\textup{(a)}] $f'(t) \le f(t)$ for every $t \in [0,1]$;
\item[\textup{(b)}] $f'(t) = f(t)$ whenever $t \notin \bigcup_l p_l$;
\item[\textup{(c)}] $\sup f' \le m_0 - 1/(4N)$.
\end{enumerate}
\end{corollary}

\begin{remark}[Why (a)+(b) are stated]
\label{rem:anchor-f-prime}
The bound (c) alone is trivialized by any constant
$f' \equiv c \le m_0 - 1/(4N)$. Condition (a) ties $f'$
pointwise below $f$; condition (b) forces $f'$ to coincide with
$f$ wherever no piece modifies it. Together they anchor $f'$ to
$f$, so generic non-constant $f$ rules out constant competitors,
and (c) becomes a non-trivial quantitative assertion about the
specific competitor the proof produces
(concretely, $f'(t) = f(t) - \sum_{l :\, t \in p_l} s_l$).
\end{remark}

The bookkeeping underlying Corollary~\ref{thm:intro-core} is
strictly more general than the one-parameter statement it
certifies: the same arithmetic applies to any compact, nonempty
topological space $K$ with continuous $f : K \to \mathbb{R}$ and
any scalars $(\delta, \varepsilon)$ with
$0 < \varepsilon \le \delta$, taking as input any finite cover of
the $\delta$-super-level set of $f$ with per-piece savings $\ge
\varepsilon$ and at-most-two overlap. The Lean library ships
both forms (Corollaries~\ref{cor:core-sup}~and~\ref{thm:main} in
\S\ref{sec:app-spec}); Theorem~\ref{thm:intro-scalar} is the
combinatorial construction that turns local-witness data into
such a cover in the 1D case.

Both results reappear in §\ref{sec:spec}:
Theorem~\ref{thm:intro-scalar} as the library's main theorem
(Theorem~\ref{thm:core},~\S\ref{sec:core-spec}), and
Corollary~\ref{thm:intro-core} as the one-parameter sup-reduction
corollary (Corollary~\ref{thm:main},~\S\ref{sec:app-spec}).

\textbf{Scope.} This note formalizes the proof of the Claim
(the combinatorial argument on $[0,1]$) as a Lean 4 library,
\CombArg. We do \emph{not}
formalize \cite[Lemma~3.1]{DLT13} (the replacement family), the
$(U^1, U^2) \in \mathcal{CO}$ disjoint-pair structure, or the
final contradiction against $m_0 = \inf_\Lambda \mathcal{F}$;
these depend on the GMT and admissible-class framework of the
ambient construction and remain the consumer's responsibility.
The library ships one main theorem (the combinatorial core,
§\ref{sec:core-spec}) on $[0,1]$, together with two
sup-reduction corollaries (§\ref{sec:app-spec}): a generic
compact-$K$ form and its one-parameter specialization. The
formalization is machine-verified against \Mathlib, declares no
custom axioms, and contains no \texttt{sorry}. Source:
\url{https://github.com/MathNetwork/comb_arg}.

